\documentclass[10pt,a4paper]{amsart}
\usepackage[margin=19mm]{geometry}
\usepackage{amsmath,amssymb,mathtools}
\usepackage[scr=rsfs,cal=euler]{mathalfa}
\usepackage{xfrac}
\usepackage[unicode,hidelinks,bookmarks=false]{hyperref}
\newcommand{\ie}{i.e.\ }
\newcommand{\cf}{cf.\ }
\newcommand{\resp}{respectively\ }
\newcommand{\Subsection}{\subsection}
\providecommand{\nobreakdash}{\nobreak\hbox{-}}
\setlength{\emergencystretch}{2em}
\multlinegap0pt
\usepackage{bm}
\usepackage{bbm}
\usepackage{dsfont}
\usepackage{xspace}
\usepackage{cancel}
\usepackage{mathtools}
\usepackage{amsthm}
\makeatletter
\@ifundefined{thm}{\newtheorem{thm}{Thm}}{}
\@ifundefined{lem}{\newtheorem{lem}[thm]{Lemma}}{}
\@ifundefined{prop}{\newtheorem{prop}[thm]{Proposition}}{}
\makeatother
\newcommand{\Bdrep}{\mathrm{B}^\mathrm{drep}}
\newcommand{\Ndrep}{\mathrm{N}^\mathrm{drep}}
\newcommand{\Nsigma}{\mathrm{N}^\sigma}
\newcommand{\Sphere}{\mathbb{S}}
\newcommand{\Ztwo}{{\mathbb{Z}/2}}
\newcommand{\gp}{\mathrm{gp}}
\newcommand{\rE}{\mathrm{E}}
\newcommand{\sd}{\mathrm{sd}}
\title[Motivic real topological Hochschild spectrum]{Motivic real topological Hochschild spectrum}
\author{Doosung Park}
\date{Source: arXiv:2305.04150v3 (CC BY 4.0)}
\begin{document}
\maketitle

\noindent\emph{Source and licence.} Motivic real topological Hochschild spectrum. Version arXiv:2305.04150v3 (CC BY 4.0), distributed under the Creative Commons Attribution 4.0 International licence, as recorded in the arXiv licence metadata for this exact version and shown on the abstract page https://arxiv.org/abs/2305.04150v3. Full rights notice: licenses/arxiv-2305.04150-CC-BY-4.0.txt.\par\smallskip

\noindent\textit{Editorial scope.} Counted unit: the Prop whose source label is \texttt{dih.15} in the article (source0.tex lines 1143--1152, source label \texttt{dih.15}), delivered together with the article's own complete original proof of it (source0.tex lines 1153--1218, one \texttt{proof} environment of 66 lines). No mathematics is written, repaired or re-derived: statement and proof are the article's own text line for line. Required context retained uncounted: drep.6 (lines 976--986); dih.25 (lines 1103--1112). Every definition, display and label of those lines is retained unchanged. Omitted from this package and not redistributed: 1--975, 987--1102, 1113--1142, 1219--3824. Nothing inside a retained excerpt is omitted or altered: every retained body is delivered exactly as the article's lines are written, so no omission receipt of any kind accompanies this package. Citations: the retained excerpts carry no citation command at all, so no bibliography entry of the article is transcribed and no reference is redistributed. Reference adaptation: none was needed and none was made; every reference inside the delivered unit resolves inside it, so no unresolved reference or citation is produced. Printed slips kept literal: no printed word, symbol, hypothesis or number of the counted statement or of its proof was altered. Layout and class: the article's own class is \texttt{amsart} with packages including xcolor, hyperref, todonotes, amsmath, amsthm, amssymb, mathrsfs, tikz-cd; the wrapper is the lane's pinned \texttt{amsart} editorial wrapper at 10pt on A4, which restates the theorem-like environments the retained text needs and adds the article's own macro definitions that the retained text needs (\texttt{\textbackslash Bdrep}, \texttt{\textbackslash Ndrep}, \texttt{\textbackslash Nsigma}, \texttt{\textbackslash Sphere}, \texttt{\textbackslash Ztwo}, \texttt{\textbackslash gp}, \texttt{\textbackslash rE}, \texttt{\textbackslash sd}), with extra wrapper packages (bm, bbm, dsfont, xspace, cancel, mathtools). The delivered numbering is the unit's own. Assets: no retained span contains a figure environment, a graphics-inclusion command, a PDF-page inclusion, a table or a TikZ picture; no image file is copied into this package, the package declares no asset dependency and no image file is redistributed with it. Licence: version arXiv:2305.04150v3 of this article is distributed under the Creative Commons Attribution 4.0 International licence, as recorded in the arXiv licence metadata for this exact version and shown on the abstract page https://arxiv.org/abs/2305.04150v3. A full rights notice is delivered at licenses/arxiv-2305.04150-CC-BY-4.0.txt. Characters: every retained excerpt is pure ASCII, so the delivered engine, XeLaTeX with its default Latin Modern text font, reports no missing character.
\begin{lem}
\label{drep.6}
Let $P\to Q$ be a map of commutative monoids with involutions.
Then there is a natural equivalence of $\Ztwo$-spectra
\[
\Sphere[P\amalg P]\wedge_{\Sphere[P]}\Sphere[Q]
\simeq
\Sphere[(P\amalg P)\oplus_P Q],
\]
where the involution on $P\amalg P$ in the formulation switches the components.
\end{lem}

\begin{prop}
\label{dih.25}
Let $P$ be a commutative monoid with involution.
Then there is a natural isomorphism of dihedral sets
\[
\Ndrep P
\simeq
P\times \Nsigma P^\gp.
\]
\end{prop}

\begin{prop}
\label{dih.15}
Let $P$ be a commutative monoid with involution.
Then there is a natural equivalence of $\Ztwo$-spectra
\[
\Sphere[\Bdrep P]
\simeq
\Sphere[P]\wedge_{\Sphere[(i_\sharp i^*P)^{ex}]} \Sphere[P].
\]
\end{prop}

\begin{proof}
Let $\rE P^\gp$ denote the total simplicial set of $P^\gp$,
and let $Q$ be the real simplicial set whose underlying simplicial set is $P\times \rE P^\gp$,
and whose involution is given by
\[
(x,g_0,\ldots,g_q)\in P\times (P^\gp)^{\times(q+1)}
\mapsto
(w(x),w(x)-w(g_q),\ldots,w(x)-w(g_0))
\]
in simplicial degree $q$ for every integer $q\geq 0$.
Consider the maps of real simplicial sets
\[
(i_\sharp i^* P)^{ex} \to Q \to P\times \Nsigma P^\gp
\]
given by
$
(x,g)
\mapsto
(x,g,\ldots,g)
$
and
$
(x,g_0,\ldots,g_q)
\mapsto
(x,g_1-g_0,\ldots,g_q-g_{q-1})
$
in simplicial degree $q$.

There is a real simplicial isomorphism
\[
Q \oplus_{(i_\sharp i^*P)^{ex}} P
\simeq
P\times \Nsigma P^\gp.
\]
Since $Q$ is degreewise the disjoint union of finite copies of $(i_\sharp i^*P)^{ex}$ and $(i_\sharp i^*P)^{ex}\amalg (i_\sharp i^*P)^{ex}$ with the switching involution as an $(i_\sharp i^*P)^{ex}$-set, Lemma \ref{drep.6} yields a natural equivalence of $\Ztwo$-spectra
\begin{equation}
\label{dih.15.1}
\Sphere[Q] \wedge_{\Sphere[(i_\sharp i^*P)^{ex}]} \Sphere[P]
\simeq
\Sphere[P\times \Nsigma P^\gp].
\end{equation}

Let us show that the map $Q\to P$ given by
\[
(x,g_0,\ldots,g_q)\mapsto x
\]
in simplicial degree $q$ is a $\Ztwo$-homotopy equivalence. Its underlying map of simplicial sets is the projection $P\times \rE P^\gp \to P$, which is a homotopy equivalence.
The $\Ztwo$-fixed point of the Segal subdivision $\sd_{\sigma} Q$ is in simplicial degree $q$ the set
\[
\{
(x,g_0,\ldots,g_q,w(x)-w(g_q),\ldots,w(x)-w(g_0))
:
x\in P^{\Ztwo},g_0,\ldots,g_q\in P^\gp
\}.
\]
From this description,
we see that the induced map
\[
(\sd_{\sigma} Q)^{\Ztwo}
\to
(\sd_{\sigma} P)^{\Ztwo}
\]
can be identified with the projection $P\times \rE P^\gp \to P$.
Hence the map $Q\to P$ is a $\Ztwo$-homotopy equivalence.
Combine this with Proposition \ref{dih.25} and \eqref{dih.15.1} to obtain the desired equivalence.
\end{proof}

\end{document}
